\chapter{The European Cloud Landscape}
\label{ch:landscape}

Before any provider could be measured, there had to be a list of providers worth
measuring. Europe has hundreds of companies that rent out servers in one form or
another, and most of them could never take over the work Alps Alpine runs on
\acs{AWS}. This chapter explains how the eighteen providers in this thesis were
found and chosen, and then introduces each of them briefly, so that the names that
fill the following chapters are familiar before the studies begin.

\section{How the Providers Were Chosen}
\label{sec:provider-selection}

Most of the providers came from European Alternatives, a website that collects
European replacements for digital services and products, including cloud
services \cite{EuropeanAlternatives_Cloud}. Ten of the providers in this thesis
are listed there: Hetzner, IONOS, OVHcloud, Scaleway, Elastx, Exoscale,
Infomaniak, Cyso Cloud, UpCloud and Deutsche Telekom's public cloud, which the
site lists as T~Cloud \cite{EuropeanAlternatives_Cloud,EuropeanAlternatives_VPS}.
Deutsche Telekom's platform was also chosen because of the large community of
users behind it.

Two providers came from within Alps Alpine itself \cite{AlpsAlpine_Internal_2026}.
STACKIT was recommended internally as a candidate worth examining, and Arvato
Systems was included because it had approached the company in the past.

The remaining providers were added because of their part in Europe's own efforts
to build sovereign cloud services. plusserver is a founding member of Gaia-X, the
European initiative described in \cref{sec:cloud-models}, and helped start the
Sovereign Cloud Stack, an open-source cloud platform developed in Germany
\cite{Plusserver_Associations}. Cleura, then called City Network, joined the
Gaia-X project in 2020 \cite{Cleura_GaiaX_2020}. noris network offers a sovereign
cloud certified under the Sovereign Cloud Stack \cite{SCS_noris}. Nine was added
as a Swiss provider that keeps its customers' data in certified data centres in
Switzerland \cite{Nine_Infrastructure}. Finally, a locally hosted setup, in which
the company runs its own servers, was included as a point of reference rather than
as a provider.

Being on one of these lists was only the first step. To stay in the study, a
provider also had to meet six conditions:

\begin{enumerate}
  \item it is registered in a member state of the \ac{EU} or the \ac{EEA}, since
        the law that applies to a provider follows where it is registered;
  \item it sells cloud infrastructure that a customer can set up and control
        itself, not only hosting, rented rack space or finished software;
  \item it runs its data centres in Europe and does not, by default, copy or route
        data elsewhere;
  \item it runs its own platform rather than reselling one of the American
        providers', because a reseller simply inherits the position of the
        platform beneath it;
  \item it publishes a list of its services detailed enough to compare with the
        services Alps Alpine uses; and
  \item it publishes enough about its technology and governance for the
        sovereignty questions to be answered from public sources.
\end{enumerate}

The last condition has a consequence worth stating at the outset. A provider that
publishes little cannot be assessed fully, so a low score in this thesis may
reflect poor documentation as much as poor practice
(\cref{sec:design-limitations}).

\Cref{tab:providers} lists the providers chosen and how each was found. They come
from six countries and include every common kind of owner: independent firms,
companies listed on the stock market, subsidiaries of telecoms groups and a
retail group. Deutsche Telekom's platform appears twice, under its old name and its
new one, for reasons explained in \cref{sec:study1}.

\begin{table}[htbp]
  \centering
  \caption{The eighteen assessed providers and how each was found.}
  \label{tab:providers}
  \footnotesize
  \begin{tabular}{@{}l l c >{\raggedright\arraybackslash}p{4.3cm}
                           >{\raggedright\arraybackslash}p{3.5cm}@{}}
    \toprule
    \textbf{Provider} & \textbf{Country} & \textbf{Founded} & \textbf{Ownership}
      & \textbf{Found through} \\
    \midrule
    Hetzner            & Germany     & 1997 & Independent                & European Alternatives \\
    IONOS              & Germany     & 1988 & United Internet AG         & European Alternatives \\
    STACKIT            & Germany     & 2018 & Schwarz Group              & Internal communication \\
    T-Systems \acs{OTC} & Germany     & 2016 & Deutsche Telekom AG        & European Alternatives \\
    T-Cloud Public     & Germany     & 2026\textsuperscript{a} & Deutsche Telekom AG        & European Alternatives \\
    Arvato Systems     & Germany     & ---  & Bertelsmann SE \& Co.\ KGaA & Internal communication \\
    noris network      & Germany     & 1993 & Independent                & Sovereign Cloud Stack \\
    PlusServer         & Germany     & 1999 & BC Partners funds\textsuperscript{c} & Gaia-X \\
    OVHcloud           & France      & 1999 & Listed; Klaba family       & European Alternatives \\
    Scaleway           & France      & 1999 & Iliad Group                & European Alternatives \\
    Cleura             & Sweden      & 2009\textsuperscript{b} & Iver (EQT)                 & Gaia-X \\
    Elastx             & Sweden      & 2012 & Privately held             & European Alternatives \\
    Exoscale           & Switzerland & 2011 & A1 Telekom Austria Group   & European Alternatives \\
    Infomaniak         & Switzerland & 1994 & Independent; Infomaniak Foundation & European Alternatives \\
    Nine               & Switzerland & 2000 & Independent                & Swiss data residency \\
    Cyso Cloud         & Netherlands & 2016 & Cyso Group                 & European Alternatives \\
    UpCloud            & Finland     & 2011 & Privately held             & European Alternatives \\
    Locally hosted     & ---         & ---  & The organisation itself    & Reference point \\
    \bottomrule
  \end{tabular}

  \vspace{2pt}
  {\footnotesize\raggedright \textsuperscript{a}~The new name of the Open Telekom
  Cloud; the platform itself dates from 2016. \textsuperscript{b}~Launch of its
  cloud; the company, then City Network, dates from the early 2000s.
  \textsuperscript{c}~Last documented owner, since 2017.\par}
\end{table}

\section{The Providers}
\label{sec:provider-profiles}

\subsection{Germany}

\subsubsection*{Hetzner}
Hetzner is an independent German company founded in 1997 in Gunzenhausen,
Bavaria, and best known for inexpensive servers \cite{Hetzner_About}. Its cloud runs from Germany and Finland,
and also from the United States and Singapore \cite{Hetzner_Locations}. It
represents the independent, low-cost end of the German market.

\subsubsection*{IONOS}
IONOS grew out of 1\&1, a German company founded in 1988, and is part of the
United Internet group \cite{IONOS_About}. Long known for websites and domains, it also sells
cloud infrastructure, and represents the established German hosting market in this
thesis.

\subsubsection*{STACKIT}
STACKIT is the cloud of the Schwarz Group, the German retail group. It was
developed from 2018 for the group's own use and opened to outside customers in 2022
\cite{STACKIT_About,STACKIT_Launch_2022}.
Of the European providers studied in depth, it is the most recent to open to
customers and the one with the largest owner behind it (\cref{sec:method-study3}).

\subsubsection*{T-Systems Open Telekom Cloud and T-Cloud Public}
Deutsche Telekom launched the Open Telekom Cloud in 2016, built with hardware and
software from Huawei and run from its data centres in Magdeburg and Biere
\cite{OTC_GoesLive_2016,TCloudPublic_SecureCloud}. At the start of 2026 it
renamed the platform T-Cloud Public \cite{TCloudPublic_2026}. Both names are assessed, because
the difference between them turned out to be a finding in its own right.

\subsubsection*{Arvato Systems}
Arvato Systems, based in G\"utersloh, is a wholly owned subsidiary of
Bertelsmann, the German media and services group. It runs its own data centres in
Germany and offers managed cloud services mainly to large companies, alongside a
private cloud of its own \cite{Arvato_About,Arvato_DataCenter}.

\subsubsection*{noris network}
noris network is an independent, owner-managed data-centre operator from
Nuremberg, founded in 1993, which runs six data centres in Germany
\cite{NorisNetwork_About}. Its sovereign cloud is built on, and certified under,
the Sovereign Cloud Stack \cite{SCS_noris}.

\subsubsection*{plusserver}
plusserver is a provider from Cologne, bought in 2017 by funds advised by the
investment firm BC Partners \cite{PlusServer_About,BCPartners_Plusserver_2017}. It
runs its cloud from four data centres in Hamburg, Cologne and D\"usseldorf
\cite{Plusserver_DataCenters}, is a founding member of Gaia-X and was one of the
initiators of the Sovereign Cloud Stack \cite{Plusserver_Associations}.

\subsection{France}

\subsubsection*{OVHcloud}
OVHcloud was founded in Roubaix in 1999 by Octave Klaba, whose family still holds
most of its shares, and has been listed on the Paris stock exchange since 2021
\cite{OVHcloud_IPO,OVHcloud_Shareholders}. It is the
only provider in the study that builds its own servers, at a factory in northern
France \cite{NextInk_OVHCroix_2023}.

\subsubsection*{Scaleway}
Scaleway, which traces its origins to the French hosting company Online in 1999,
belongs to the French telecoms group iliad. It runs its cloud from Paris,
Amsterdam, Warsaw and Milan \cite{Scaleway_About}, is aimed mainly at software
developers, and offers a growing range of ready-made services.

\subsection{Sweden}

\subsubsection*{Cleura}
Cleura, known until 2022 as City Network, is a Swedish provider based in
Karlskrona that launched one of Sweden's first cloud services in 2009. Since 2020
it has belonged to the Swedish IT group Iver, itself owned by the investment firm
EQT \cite{Cleura_About,Cleura_Iver_2020}. It took part in Gaia-X from its early
days \cite{Cleura_GaiaX_2020}.

\subsubsection*{Elastx}
Elastx is a small, privately held Swedish provider founded in Stockholm in 2012.
It runs three data centres in the Stockholm region \cite{Elastx_About}, and
represents the smaller end of the market.

\subsection{Switzerland}

\subsubsection*{Exoscale}
Exoscale was founded in Lausanne in 2011 and now belongs to A1 Digital, part of
the A1 Telekom Austria Group. Its cloud runs from Switzerland, Austria, Germany,
Bulgaria and Croatia \cite{Exoscale_About}. It brings a telecoms owner from outside
Germany and France into the study.

\subsubsection*{Infomaniak}
Infomaniak was founded in Geneva in 1994 and runs its data centres in Geneva. It
is independent, and a foundation now holds the shareholder rights for the purpose
of keeping it that way \cite{Infomaniak_About}.

\subsubsection*{Nine}
Nine is an independent Swiss provider based in Zurich, which has offered hosting
and managed services in Switzerland since 2000 \cite{Nine_About}. It runs its
services from two certified data centres in the Zurich area
\cite{Nine_Infrastructure}.

\subsection{The Netherlands and Finland}

\subsubsection*{Cyso Cloud}
Cyso Cloud is a Dutch cloud built on the open-source software OpenStack. It began
in 2016 as Fuga Cloud and took its present name in 2025, when Cyso, a Dutch hosting
company founded in 1997 and part of the same group, took it over. It runs from
Amsterdam and Frankfurt \cite{Cyso_FugaComesOfAge,FugaCloud_About}.

\subsubsection*{UpCloud}
UpCloud is a privately held Finnish provider from Helsinki \cite{UpCloud_About}.
Alongside its European data centres it also runs locations in the United States,
Singapore and Australia \cite{UpCloud_DataCenters}.

\subsection{Locally Hosted Infrastructure}

The final entry is not a company but a situation: an organisation running its own
servers, with its own staff, in buildings it controls. It is included because it
marks one end of the range. It offers the most control a company can have and, as
\cref{sec:study2} shows, very few ready-made services.

\section{AWS as the Point of Comparison}
\label{sec:aws-reference}

\acs{AWS} is not one of the eighteen, but it appears throughout the thesis as the
yardstick against which the others are measured. Alps Alpine already runs its
work there, so the practical question is how much of that work each European
provider could take over, and on what terms. \acs{AWS} is therefore used as the
reference in the capability, financial and benchmarking studies, but it is not
scored for sovereignty, since the laws described in \cref{subsec:cloud-act} are
the reason the search for alternatives began.
